\chapter{Urinary Tract Infection}
\index{Urinary Tract Infection}

\medskip
\noindent\textit{\textbf{Under review.} This chapter is an AI-assisted educational draft; it carries no source citations and is not a substitute for professional medical advice.}

\section*{Definition and Overview}
A urinary tract infection (UTI) is an infection of any part of the urinary system, most commonly the bladder (cystitis) and, when it travels higher, the kidneys (pyelonephritis). It is usually caused by bacteria entering the urinary tract through the urethra and multiplying. Most UTIs are simple, uncomplicated infections that settle well with antibiotics, but they can recur and, if untreated, spread to the bloodstream.

\section*{Epidemiology}
UTIs are among the most common infections seen in general practice. They are much more common in women, and roughly 1 in 3 women will experience at least one by mid-life. Risk increases with sexual activity, pregnancy, menopause, urinary catheters, and conditions that cause incomplete bladder emptying. Men have fewer UTIs, but when they occur they may reflect prostate enlargement. Recurrent infections affect a significant minority of women.

\section*{Aetiology and Pathophysiology}
The great majority of UTIs are caused by bacteria that normally live in the bowel, most often \textit{Escherichia coli}. Factors that allow bacteria to reach and multiply in the bladder include:

\begin{itemize}
    \item \textbf{Bowel bacteria} (e.g.\ \textit{E. coli}) ascending the urethra -- the most common mechanism
    \item \textbf{The short urethra in women}, which makes ascent easier
    \item \textbf{Sexual activity}, which can push bacteria into the urethra
    \item \textbf{Incomplete bladder emptying} from prostate enlargement, prolapse, or nerve conditions
    \item \textbf{Urinary catheters} or instruments in the bladder
    \item \textbf{Pregnancy}, which slows the flow of urine
    \item \textbf{Menopause}, with changes in the vaginal environment
    \item \textbf{Diabetes} and other causes of a weakened immune system
\end{itemize}

Once in the bladder, bacteria attach to the lining and multiply, causing inflammation. If the infection is not cleared, it can travel up the ureters into the kidneys.

\section*{Clinical Features}
Bladder infection (cystitis) is the most common presentation:
\begin{itemize}
    \item Pain or burning when passing urine (dysuria)
    \item Passing urine frequently and urgently
    \item A feeling that the bladder is not empty
    \item Cloudy, strong-smelling, or blood-stained urine
    \item Lower abdominal or suprapubic discomfort
\end{itemize}

Kidney infection (pyelonephritis) causes more general symptoms:
\begin{itemize}
    \item Fever, chills, and feeling generally unwell
    \item Flank or loin pain (back pain on one or both sides)
    \item Nausea and vomiting
    \item Symptoms of cystitis are often present as well
\end{itemize}

\section*{Differential Diagnosis}
\begin{itemize}
    \item Urethritis from sexually transmitted infections (e.g.\ chlamydia, gonorrhoea)
    \item Bladder stones or foreign bodies
    \item Bladder cancer (blood in the urine)
    \item Kidney stones causing pain
    \item Interstitial cystitis (bladder pain without infection)
    \item Vaginal infections causing irritation
    \item Prostatitis in men
\end{itemize}

\section*{When to Seek Medical Care}
Most uncomplicated UTIs in women can be managed by a GP. Seek assessment if symptoms are severe, not improving after a day or two, or accompanied by fever or flank pain.

\textbf{Contact your local emergency services for:}
\begin{itemize}
    \item High fever with shaking chills or confusion
    \item Inability to pass any urine
    \item Signs of sepsis: rapid breathing, rapid heart rate, mottled skin, or low blood pressure
    \item Severe flank or abdominal pain with vomiting
\end{itemize}

\section*{Diagnosis and Investigations}
\begin{itemize}
    \item \textbf{History and examination:} symptoms, prior UTIs, sexual and gynaecological history in women, prostate symptoms in men
    \item \textbf{Urine dipstick:} detects white blood cells and nitrites that suggest infection
    \item \textbf{Urine culture:} identifies the bacteria and its antibiotic sensitivities; used when infection is complicated or recurrent
    \item \textbf{Blood tests:} in severe cases, to check kidney function and markers of infection
    \item \textbf{Imaging:} ultrasound or CT for recurrent infections or suspected structural problems
    \item \textbf{Referral:} to urology for men, children, or persistent cases
\end{itemize}

\section*{Management}
\begin{itemize}
    \item \textbf{Antibiotics:} a short course is usually prescribed for bacterial cystitis; a longer or intravenous course for kidney infection
    \item \textbf{Pain relief:} paracetamol or ibuprofen
    \item \textbf{Fluids:} drink normally unless a clinician has advised a restriction; forcing large volumes does not replace treatment
    \item \textbf{Symptom relief:} urinary alkalinisers (e.g.\ sodium citrate) may ease burning
    \item \textbf{Delayed antibiotic approach:} for mild cystitis, some clinicians offer a prescription to use only if symptoms do not settle
    \item \textbf{Recurrent infections:} may be managed with prevention strategies or, in some cases, a low-dose antibiotic plan
    \item \textbf{Admission to hospital:} for kidney infection with sepsis, vomiting, or pregnancy complications
\end{itemize}

\section*{Complications}
\begin{itemize}
    \item Kidney infection (pyelonephritis)
    \item Abscess within or around the kidney
    \item Sepsis, a life-threatening whole-body response
    \item Recurrent infections
    \item Kidney damage in recurrent or severe infection
    \item Complications in pregnancy, including preterm birth and low birth weight
    \item Antibiotic resistance
\end{itemize}

\section*{Prognosis and Outcomes}
Uncomplicated cystitis almost always settles within a few days of treatment. Kidney infections improve more slowly and may take a week or two to fully resolve. Recurrent infections can be reduced with prevention measures. Severe infections are more dangerous in older people, in those with diabetes or catheters, and in pregnant women.

\section*{Prevention and Self-Care}
\begin{itemize}
    \item Drink enough to avoid dehydration, following individualized advice if you have kidney, heart, or fluid-balance problems
    \item Pass urine after sexual intercourse
    \item Wipe from front to back after using the toilet
    \item Do not delay urination when you feel the urge
    \item Avoid perfumed soaps, sprays, and other products that irritate the area
    \item Wear cotton underwear and avoid tight, synthetic clothing
    \item Some people may benefit from cranberry products or vaginal oestrogen, on medical advice
    \item Complete prescribed antibiotic courses as directed
\end{itemize}

\section*{Sources and Further Reading}
\begin{itemize}
    \item NHS. \textit{Urinary tract infections (UTIs).} https://www.nhs.uk/conditions/urinary-tract-infections-utis/
    \item National Institute for Health and Care Excellence. \textit{Urinary tract infection guidance.} https://www.nice.org.uk/guidance
\end{itemize}
